\section{Related Work}
\label{sec:related-work}

\textbf{Recommendation interfaces and augmentation.} Language-based recommendation turns histories, item content, and targets into interfaces learnable by pretrained language models. P5~\citep{geng2022p5} formulates recommendation as personalized text generation, TALLRec~\citep{bao2023tallrec} aligns an LLM through recommendation tuning, and TIGER~\citep{rajput2023tiger} and GRLM~\citep{zhang2026grlm} generate semantic or textual item identifiers. LLaRA~\citep{liao2024llara} and E4SRec~\citep{li2023e4srec} instead connect collaborative or ID-based signals with LLM representations. Sequential recommendation provides a broad augmentation vocabulary~\citep{dang2024survey}: CL4SRec~\citep{xie2022cl4rec} and CoSeRec~\citep{liu2021coserec} transform interaction sequences, DuoRec~\citep{qiu2022duorec} and ICLRec~\citep{chen2022iclrec} construct semantic or intent-aware contrastive pairs, and EC4SRec~\citep{wang2022ec4srec} uses explanation-guided perturbations. Recent generative methods sample histories and targets in GenPAS~\citep{lee2025genpas}, learn verbalizations in From Logs to Language~\citep{shi2026verbalization}, augment conventional recommendation in Llama4Rec~\citep{luo2024llama4rec}, and use LLMs as data enhancers in LLMSeR~\citep{wang2025llmser}. These methods change the available supervision, but applying one method uniformly fixes its user population and exposure implicitly. Our experiments include all baseline methods listed in Table~\ref{tab:overall}, with L2Aug~\citep{wang2022l2aug}, GenPAS~\citep{lee2025genpas}, and AsarRec~\citep{zhang2026asarrec} as direct recommendation-augmentation comparators.

\textbf{Adaptive policy search and data curation.} Automated augmentation searches transformation compositions and strengths through AutoAugment~\citep{cubuk2019autoaugment}, Population Based Augmentation~\citep{ho2019pba}, and RandAugment~\citep{cubuk2020randaugment}. Recommendation methods introduce finer conditioning: CoSeRec~\citep{liu2021coserec} restricts transformations by sequence length, L2Aug~\citep{wang2022l2aug} learns an editing policy from recommender feedback, and AsarRec~\citep{zhang2026asarrec} learns user-conditioned transformation matrices. In parallel, data-allocation methods optimize which sources receive training opportunity: DoReMi~\citep{xie2023doremi} learns domain weights, Skill-It!~\citep{chen2023skillit} schedules skills from dependency estimates, and Data Mixing Laws~\citep{ye2024mixinglaws} and RegMix~\citep{liu2024regmix} predict favorable mixtures from smaller runs. DataComp~\citep{gadre2023datacomp}, Data-Juicer~\citep{chen2024datajuicer}, and JEST~\citep{evans2024jest} systematize corpus evaluation, processing, and learnability-aware selection. These lines establish validation-guided policy search and exposure control, but their action spaces typically remain within one transformation family or operate over domains, examples, or corpora. Our direct comparison includes L2Aug~\citep{wang2022l2aug}, GenPAS~\citep{lee2025genpas}, and AsarRec~\citep{zhang2026asarrec} as recommendation-augmentation baseline methods.

\textbf{Agents for training-data design.} Language agents provide the control loop needed to operationalize such data decisions. ReAct~\citep{yao2023react} couples reasoning with tool actions, Reflexion~\citep{shinn2023reflexion} carries feedback into later attempts, and ADAS~\citep{hu2025adas} and AgentSquare~\citep{shang2025agentsquare} search structured agent designs. Curation-Bench~\citep{kang2026curation} applies an inspect--implement--train--evaluate loop directly to data policies and shows how downstream outcomes can revise later proposals. Recommendation augmentation adds execution constraints absent from general free-form curation: assignments must identify users, respect method applicability and exposure, preserve required task fields, and retain provenance. These agent methods motivate our controller; the recommendation-specific evaluation includes all baseline methods listed in Table~\ref{tab:overall}.
